
\paragraph{Uncertainty, reserves and sub-hourly operation.}
The dispatch is deterministic and hourly. Price and demand are perfect-foresight time
series; forecast error, reserve provision, intra-hour demand excursions and operational
disturbances such as unit outages are all outside the formulation. Each of these plausibly
raises the value of the physics this paper finds inert. Transport delay is a zero-timestep
shift on an hourly grid and would not be at a five-minute one; thermal inertia in the pipe
volume buffers exactly the sub-hourly excursions excluded here; and storage and heat-pump
flexibility are worth more when they can be sold into reserve markets than when they can
only arbitrage a known price curve. The results should therefore be read as establishing
that extended network physics is decision-irrelevant \emph{in the deterministic hourly
planning regime}, the regime in which the great majority of published district-heating
dispatch models operate, and not as a claim about real-time control. The direction of
the bias is worth stating plainly: relaxing these assumptions would generally give the
extended physics more room to matter, though not necessarily monotonically, so the nulls
reported here are best read as conservative within the tested regime; the loss-dominance
algebraic identity does not depend on them, though its empirical magnitude is conditional on
the tested systems.

A related caution applies to the estimation-bias result in the opposite direction.
\citet{Koek2026} show for multistage investment planning that ignoring price and
demand-trajectory uncertainty, the value of the stochastic solution, biases technology
selection by 27--39\,\%. Their setting is investment rather than dispatch and does not
vary network topology, but the comparison suggests that the resolution bias quantified
here and the uncertainty bias they quantify are additive rather than substitutable. A
model can be wrong about the network and wrong about the future independently, and
correcting one does not correct the other; addressing both within a single formulation
remains open.

\paragraph{Demand-side flexibility.}
All demands are exogenous time series and demand response is excluded. This is one of the
few exclusions that could plausibly \emph{increase} the topology effect rather than leave
it unchanged, because spatially distributed flexibility has location-dependent value: a
flexible load near the network far end can absorb local surplus without incurring the
trunk losses that serving it from the plant would. Whether that is enough to lift routing
above the decision threshold is untested here and is a well-posed question for a follow-up
study.

\paragraph{Network topology class and generation siting.}
Every case, real and synthetic, is a radial tree with unidirectional flow. Ring and meshed
topologies are not merely a larger instance of the same problem: with pipe-flow direction
fixed by the pipe list, the return-side pressure balance around a cycle forces the sum of
the loop's friction drops to be non-positive, which is infeasible for any non-zero flow.
Admitting loops requires signed-flow hydraulic variables, a per-pipe direction selection
that reorients each pressure drop and temperature mix, which is a distinct formulation
rather than a parameter change. The same boundary limits the generation-topology
moderator: the synthetic generator roots all injection at the primary producer, so the
distributed-generation case is posed as a controlled open question rather than answered,
and multi-source operation on a meshed network is left to the companion study. Because
generation is central in every case examined, the null reported for spatial routing is
established for centrally-supplied networks specifically.

\paragraph{The nonlinear reference.}
The nonlinear reference is solved on representative 72-hour windows rather than globally. The
full-year non-convex programme finds no incumbent within the solver budget, and a low-load
summer 72-hour window is infeasible under native physics with the linearised schedule fixed. Both
are reported as results in Section~\ref{subsec:linearisation} (the second is a solver outcome
whose physical-versus-numerical cause is not yet established; an IIS diagnosis is left to future
work), but they do bound the claim: the solved linearisation error is established for winter
and autumn operating points, and whether it remains sub-percent under shoulder-season
operation with frequent on/off cycling is open. Warm-start heuristics or a reformulation
of the bilinear terms are the natural routes.

\paragraph{One real network, and the synthetic generator's assumptions.}
The Memmingen figures come from a single industrial network; the synthetic factorial
supplies external validity for the \emph{structure} of the result (loss dominance,
adder non-transfer, the loss-number rule) but not for the absolute magnitudes, which
should not be treated as universal constants. The synthetic networks also share the
assumptions of the generator that produced them: radial trees, a common heating curve,
proportionally sized capacity under a stated convention. They are a controlled parametric
family, not a sample of real networks, and the balanced factorial buys unambiguous
attribution across the design at the cost of that representativeness.

\paragraph{Demand-concentration metric.}
The normalised variance index used as a synthetic design factor has a narrow dynamic range
on realistic tree networks, and the model-selection observations in
Table~\ref{tab:criteria} therefore use the Gini coefficient and top-$k$ demand share
instead. An empirical mapping from either metric to the aggregation gap would require
several real networks and is a natural follow-up.
